\section*{ID9008: Relation: |p I - Π(V α)|\textasciicircum{}2}
\paragraph{Metadata.}
PiID: 9536534940603 | RTEID: RTE:P34980.0645:T1.0877 | UUID: 11e7ffad-5a92-5862-8c0f-2b0961c3f800

\paragraph{Canonical Statement.}
"Geometric Weighting Kernel," defining the degree of influence or "affinity" that each high-dimensional logic pole exerts over the localized voxels of the spherical lattice. This identifier formalizes the process of Topological Projection , wherein the abstract 4D logic of the tesseract is mapped onto the 3D physical surface of the vacuum vessel to facilitate real-time field navigation. The mathematical core of ID 018 is a Gaussian-weighted affinity matrix \$\textbackslash{}mathcal\{A\} \{i \textbackslash{}alpha\}\$. The term \$|\textbackslash{}mathbf\{p\} i - \textbackslash{}Pi(V \textbackslash{}alpha)|\textasciicircum{}2\$ represents the squared Euclidean distance between a specific node address \$i\$ on the 153,600-point lattice and the projection \$\textbackslash{}Pi\$ of a hypercubic vertex \$V \textbackslash{}alpha\$ into 3-space. The parameter \$\textbackslash{}sigma\_T\$ is the Trawinized Coherence Width , a scale-invariant constant that prevents informational overlap between adjacent vertices. By multiplying the resulting kernel by the Trawin Constant ( \$

\paragraph{Canonical Equation.}
\[
|\mathbf{p} i - \Pi(V \alpha)|^2
\]

\paragraph{Dictionary.}
Relation: |p I - Π(V α)|\textasciicircum{}2 — |p i - Π(V α)|\textasciicircum{}2

\paragraph{Encyclopedia.}
Relation: |p I - Π(V α)|\textasciicircum{}2 is a symbolic expression, written |p i - Π(V α)|\textasciicircum{}2. It introduces a symbol or relation used elsewhere in the framework. It is built from Π, α, p, V. Within the Trawin Zero Point registry it serves as one element of the framework's mathematical narrative.
